% Draft: abstract for the restructured paper. First drafted 18 Oct 2026 in the
% six moves of the structure review (reviews/structure-muse-20261017-115410/
% MUSE.md, §4), revised after the §1 muse review, then trimmed the same day at
% Brett's request ("trim the abstract and make it NBER-like (not strictly)").
% Rewritten from scratch 23 Oct 2026 (Brett: "throw out the abstract and start
% again, NBER style"): the puzzle, the test, the three results, and what they
% imply, in plain prose. Verdict as in §6 (ordinary Head preferred for noun
% phrases without a relative clause; fusion favoured, not forced, in fused
% relatives). Once/twice/thrice and the genitive are left to the body.

\begin{abstract}
Some grammars give a word two syntactic functions at once. \textit{The Cambridge Grammar of the English Language} treats \mention{some} in \mention{take some} as both the determiner and the head of its noun phrase, and \mention{rich} in \mention{the rich} as both modifier and head, a fusion of functions that \textcite{Payne2007} defend. The simpler alternative is that the word is just the head. I argue that a second function needs evidence of form that varies with it rather than with the word, and test for such evidence with the grammar's own diagnostics and new corpus data. For the modifier function there's none. For the determiner function there's one fact, that no adjective comes before a bare determinative, but the comparison is close, and on balance it doesn't justify the extra function. Fused relatives such as \mention{whoever told you that} are different: they relativize subjects, they have no infinitival forms, and in usage the case of \mention{whoever} answers to its function inside the clause. There fusion is favoured, though not forced. Keeping a category, heading a phrase, and taking a reading from context are things an ordinary head can do, not evidence of a second function.
\end{abstract}
